\scanpage{93}
\deleted{\unsure{Soit}} Soit $T$ l'espèce de structure ``anneau'' \inserted{(commutatif avec unité)}, $T^{\unsure{\alpha|a|\circ}}$ l'espèce de structure ``corps''. Cette dernière n'est pas définissable par $\varprojlim$ (car la catégorie des corps \inserted{(commutatifs)} n'est pas stable par $\varprojlim$, pas $\overline{\mathrm m}$ par produits finis). Leur sous-théorie pleine \inserted{de $T$} la plus petite qui soit définissable par $\varprojlim$ est \deleted{\unsure{ainsi que|admettent par} : les propriétés suivantes : sont}\inserted{définie par la sous-catégorie $\mathscr C'$ de $\mathscr C$ définie par}
\[
\begin{aligned}
C\in\mathrm{Ob}\,\mathscr C'\ \Longleftrightarrow{}&
 \ \forall\,\unsure{\mathrm{hom.}|\mathrm{Hom.}}\ A\longrightarrow B\text{ dans }(\mathrm{Ens})\text{ tel que}\\
&X=\unsure{\mathrm{Sp}|\mathrm{Spm}|h}(B)\mathbin{\unsure{-|\to}}Y=\mathrm{Sp}(A)\text{ soit bijectif : ext.}\\
&\text{résiduelles triviales, }\mathrm{Hom}_{\mathrm{ann}}(B,C)\mathbin{-}\mathrm{Hom}_{\mathrm{ann}}(A,C)\\
&\text{est bijectif.}
\end{aligned}
\]

Un tel $C$ a la propriété que $C$ est \deleted{\unsure{abs. plat}}\inserted{pseudocorps} (Spec $C$ \inserted{séparé donc} compact, \deleted{\unsure{et}} $C$ réduit \deleted{$\Longleftrightarrow$} les \deleted{anneaux locaux} anneaux locaux de $C$ des corps) -- Les $C$ sont ceux qui transforment en \unsure{isom|isom.} l'hom évident $\mathbb Z[t]\longrightarrow\mathbb Z[t,t^{-1}]\times\mathbb Z\simeq\mathbb Z[x,y]/((xy-1)x,(xy-1)y)$ transformt $t$ en $(t,0)$ resp. en $x$. La réciproque est vraie, \inserted{et connue} -- donc la catégorie \deleted{des} $\mathscr C'$ est celle des pseudocorps.

[Mais la \deleted{notion}\inserted{structure} de corps peut elle se définir par $\varprojlim$ et $\varinjlim$, du moins dans des catégories satisfaisant des conditions d'exactitude convenables ?]

Le sous-objet $A^\ast$ (des éléments inversibles) d'un \inserted{objet} anneau d'une catégorie $\unsure{\mathscr A|\mathscr C|A}$, peut-il\deleted{\unsure{le|lle}} se définir par $\varprojlim$ (ou par $\varprojlim$ et $\varinjlim$) ? Dans le premier cas, on est ramené au cas de l'objet universel dans $R=(\mathrm{Ann})^\circ$ \inserted{$=(\unsure{\mathrm{Sch\text{-}aff.}|\mathrm{Sch\ aff.}|\mathrm{Sch\,aff}}\ \deleted{\unsure{\mathrm{f.p.}}})$} qui est l'objet $\mathbb Z[t]$.
% archive sheet number: 92
